\documentclass[12pt,a4paper]{article}
\usepackage{fontspec}
\usepackage[russian]{babel}
\usepackage{amsmath,amssymb,amsfonts,amsthm}
\usepackage{geometry}
\usepackage{microtype}

\geometry{margin=2cm}
\defaultfontfeatures{Ligatures=TeX}
\setmainfont{Times New Roman}

\newtheorem{theorem}{Теорема}[section]
\newtheorem{lemma}[theorem]{Лемма}
\newtheorem{proposition}[theorem]{Утверждение}
\theoremstyle{definition}
\newtheorem{definition}[theorem]{Определение}
\newtheorem{example}[theorem]{Пример}

\newcommand{\Int}{\operatorname{Int}}
\newcommand{\Cl}{\operatorname{Cl}}
\newcommand{\Fr}{\operatorname{Fr}}

\begin{document}

\title{\textbf{Структурированный конспект лекций по топологии и геометрии}}
\author{}
\date{}
\maketitle

\textit{Этот файл собран по материалам исходного \texttt{topologia.tex}: повторы убраны, темы разложены по лекциям, а формулы и определения приведены к единой структуре.}

\tableofcontents
\newpage

\section{Лекция 1. Общая топология}

\subsection{Множества и отображения}

Пусть $X$ --- некоторое базовое множество. Если $A$ является подмножеством множества $X$, это обозначается как $A \subset X$.

\begin{definition}[Включение подмножества]
$$A \subset X \iff \forall a \in A \implies a \in X$$
\end{definition}

\begin{definition}[Операции над подмножествами]
Пусть $A, B \subset X$. Тогда определены следующие операции:
\begin{itemize}
    \item \textbf{Разность множеств:} $A \setminus B = \{a \in A \mid a \notin B\}$
    \item \textbf{Пересечение множеств:} $A \cap B = \{a \in X \mid a \in A \text{ и } a \in B\}$
    \item \textbf{Объединение множеств:} $A \cup B = \{a \in X \mid a \in A \text{ или } a \in B\}$
\end{itemize}
\end{definition}

\begin{definition}[Семейства множеств]
Пусть $\{A_i\}_{i \in I}$ --- семейство подмножеств множества $X$, индексированное множеством $I$.
Тогда
$$\bigcup_{i \in I} A_i = \{x \in X \mid \exists i_0 \in I : x \in A_{i_0}\},$$
$$\bigcap_{i \in I} A_i = \{x \in X \mid \forall i \in I,\, x \in A_i\}.$$
\end{definition}

\begin{proposition}[Законы де Моргана и дополнения]
Для любых подмножеств выполняются соотношения:
\begin{enumerate}
    \item $A \setminus B = A \setminus (A \cap B)$.
    \item
    $$A \setminus \left(\bigcup_{i \in I} A_i\right) = \bigcap_{i \in I} (A \setminus A_i),$$
    $$A \setminus \left(\bigcap_{i \in I} A_i\right) = \bigcup_{i \in I} (A \setminus A_i).$$
    \item Если $A, B \subset X$, то:
    \begin{itemize}
        \item $A \subset B \iff (X \setminus B) \subset (X \setminus A)$.
        \item $A \subset B \iff A \cap (X \setminus B) = \emptyset$.
        \item $A \cap (X \setminus B) = A \setminus B$.
    \end{itemize}
\end{enumerate}
\end{proposition}

\begin{definition}[Типы отображений]
Пусть $f: X \to Y$ --- отображение множества $X$ в множество $Y$.
Тогда:
\begin{itemize}
    \item \textbf{Инъективным} отображение называется, если
    $$\forall x_1, x_2 \in X,\; x_1 \neq x_2 \implies f(x_1) \neq f(x_2).$$
    \item \textbf{Сюръективным} отображение называется, если
    $$\forall y \in Y,\; \exists x \in X : f(x) = y.$$
    \item \textbf{Биективным} отображение называется, если оно одновременно инъективно и сюръективно.
\end{itemize}
\end{definition}

\begin{definition}[Композиция и стандартные отображения]
\begin{itemize}
    \item \textbf{Композиция:} для $f: X \to Y$ и $g: Y \to Z$ задается как $(g \circ f)(x) = g(f(x))$.
    \item \textbf{Тождественное отображение:} $I_X: X \to X$, где $I_X(x) = x$.
    \item \textbf{Константное отображение:} $C_{y_0}: X \to Y$, где $C_{y_0}(x) = y_0$ для фиксированного $y_0 \in Y$.
\end{itemize}
\end{definition}

\begin{definition}[Образ и полный прообраз]
Пусть $A \subset X$ и $B \subset Y$.
\begin{itemize}
    \item \textbf{Полный прообраз} множества $B$: $f^{-1}(B) = \{x \in X \mid f(x) \in B\}$.
    \item \textbf{Образ} множества $A$: $f(A) = \{f(x) \in Y \mid x \in A\}$.
\end{itemize}
\end{definition}

\begin{theorem}[Свойства прообразов и образов]
Пусть $f: X \to Y$ и $g: Y \to Z$. Тогда верны следующие равенства и включения:
\begin{enumerate}
    \item $f^{-1}(B \setminus C) = f^{-1}(B) \setminus f^{-1}(C)$.
    \item $f^{-1}\left(\bigcup_{i \in I} B_i\right) = \bigcup_{i \in I} f^{-1}(B_i)$.
    \item $f^{-1}\left(\bigcap_{i \in I} B_i\right) = \bigcap_{i \in I} f^{-1}(B_i)$.
    \item $(g \circ f)^{-1}(D) = f^{-1}(g^{-1}(D))$.
    \item $f(A \setminus B) \supset f(A) \setminus f(B)$.
    \item $f\left(\bigcup_{i \in I} A_i\right) = \bigcup_{i \in I} f(A_i)$.
    \item $f\left(\bigcap_{i \in I} A_i\right) \subset \bigcap_{i \in I} f(A_i)$.
    \item $f^{-1}(f(A)) \supset A$ и $f(f^{-1}(B)) \subset B$.
\end{enumerate}
\end{theorem}

\begin{proposition}[Связь образа и прообраза]
Пусть $f: X \to Y$, $A \subset Y$ и $B \subset X$. Тогда:
\begin{enumerate}
    \item $f(f^{-1}(A)) \subset A$.
    \item $f^{-1}(f(B)) \supset B$.
    \item $B \subset f^{-1}(A) \iff f(B) \subset A$.
\end{enumerate}
\end{proposition}

\subsection{Метрические и топологические пространства}

\begin{definition}[Метрическое пространство]
Пара $(X, d)$, где $X$ --- множество, а $d: X \times X \to \mathbb{R}$ --- функция расстояния, называется \textbf{метрическим пространством}, если:
\begin{enumerate}
    \item $d(x, y) = d(y, x)$ для всех $x,y \in X$.
    \item $d(x, z) \leqslant d(x, y) + d(y, z)$ для всех $x,y,z \in X$.
    \item $d(x, y) \geqslant 0$, причем $d(x, y) = 0 \iff x = y$.
\end{enumerate}
\end{definition}

\begin{definition}[Элементы метрики]
\begin{itemize}
    \item Открытый шар: $B(a, r) = \{x \in X \mid d(a, x) < r\}$.
    \item Замкнутый шар: $\overline{B}(a, r) = \{x \in X \mid d(a, x) \leqslant r\}$.
    \item Сфера: $S(a, r) = \{x \in X \mid d(a, x) = r\}$.
\end{itemize}
\end{definition}

\begin{proposition}[О шарах и сфере]
Пусть $(X, d)$ --- метрическое пространство. Тогда:
\begin{enumerate}
    \item открытый шар $B(a, r)$ открыт в топологии, порожденной метрикой $d$;
    \item замкнутый шар $\overline{B}(a, r)$ замкнут;
    \item сфера $S(a, r)$ замкнута.
\end{enumerate}
\end{proposition}

\begin{definition}[Топологическое пространство]
Семейство подмножеств $\mathcal{T}$ множества $X$ называется \textbf{топологией}, если:
\begin{itemize}
    \item $\emptyset \in \mathcal{T}$ и $X \in \mathcal{T}$.
    \item Если $V_i \in \mathcal{T}$ для всех $i \in I$, то $\bigcup_{i \in I} V_i \in \mathcal{T}$.
    \item Если $V_1, \dots, V_n \in \mathcal{T}$, то $\bigcap_{k=1}^n V_k \in \mathcal{T}$.
\end{itemize}
Элементы $\mathcal{T}$ называются \textbf{открытыми множествами}. Множество $A \subset X$ называется \textbf{замкнутым}, если $X \setminus A \in \mathcal{T}$.
\end{definition}

\begin{proposition}[Свойства открытых и замкнутых множеств]
Пусть $(X, \mathcal{T})$ --- топологическое пространство.
\begin{enumerate}
    \item Пересечение конечного числа открытых множеств открыто.
    \item Семейство замкнутых множеств содержит $X$ и $\emptyset$, замкнуто относительно произвольных пересечений и конечных объединений.
\end{enumerate}
Следовательно, в любом непустом топологическом пространстве множества $X$ и $\emptyset$ одновременно открыты и замкнуты.
\end{proposition}

\subsection{Точки множества и непрерывность}

Пусть $(X, \mathcal{T})$ --- топологическое пространство, $A \subset X$.

\begin{definition}[Классификация точек]
\begin{enumerate}
    \item Точка $a \in A$ называется \textbf{внутренней точкой} множества $A$, если существует $V \in \mathcal{T}$ такое, что $a \in V \subset A$. Множество всех внутренних точек обозначается $\Int A$.
    \item Точка $x_0 \in X$ называется \textbf{точкой прикосновения} множества $A$, если для любой ее открытой окрестности $V$ верно $V \cap A \neq \emptyset$. Множество точек прикосновения обозначается $\Cl A$.
    \item Точка $x_0 \in X$ называется \textbf{граничной точкой} множества $A$, если она является точкой прикосновения и для $A$, и для $X \setminus A$. Множество таких точек образует границу $\Fr A$.
    \item Точка $a \in X$ называется \textbf{предельной точкой} множества $A$, если для любой окрестности $V$ точки $a$ выполняется $(V \setminus \{a\}) \cap A \neq \emptyset$.
    \item Точка $a \in A$ называется \textbf{изолированной}, если существует окрестность $V$ точки $a$ такая, что $V \cap A = \{a\}$.
\end{enumerate}
\end{definition}

\begin{theorem}[Связь внутренности, замыкания и границы]
Пусть $(X, \mathcal{T})$ --- топологическое пространство. Тогда:
\begin{enumerate}
    \item $x_0$ является точкой прикосновения множества $A$ тогда и только тогда, когда $x_0$ не является внешней точкой для $A$.
    \item $\Int A = X \setminus \Cl(X \setminus A)$.
    \item $\Int(X \setminus A) = X \setminus \Cl A$.
    \item $\Cl A = X \setminus \Int(X \setminus A)$.
    \item $\Fr A = \Cl A \setminus \Int A = \Cl A \cap \Cl(X \setminus A)$.
\end{enumerate}
\end{theorem}

\begin{theorem}[Об окрестностях и внутренних точках]
Пусть $(X, \mathcal{T})$ --- топологическое пространство.
\begin{enumerate}
    \item Открытое множество является окрестностью каждой своей точки. Наоборот, если каждая точка множества $A$ является внутренней, то $A$ открыто.
    \item Если $A$ --- окрестность точки $x$ и $A \subset B$, то $B$ --- окрестность $x$.
    \item Пересечение конечного числа окрестностей точки $x$ является окрестностью $x$.
    \item Если $c$ --- точка прикосновения множества $A$ и $c \notin A$, то $c$ --- предельная точка $A$.
    \item Если $c$ не является предельной точкой $A$, то $c$ --- изолированная точка $A$ и, следовательно, $c \in A$.
\end{enumerate}
\end{theorem}

\begin{definition}[Непрерывность отображения]
Пусть $X$ и $Y$ --- топологические пространства. Отображение $f: A \to Y$ (где $A \subset X$) называется \textbf{непрерывным в точке} $a_0 \in A$, если для любой открытой окрестности $V \subset Y$ точки $f(a_0)$ существует открытая окрестность $U \subset X$ точки $a_0$ такая, что
$$f(U \cap A) \subset V.$$
\end{definition}

\begin{definition}[Предел функции]
Точка $y_0 \in Y$ называется \textbf{пределом} функции $f(x)$ при $x \to x_0$, если для любой открытой окрестности $V$ точки $y_0$ существует открытая окрестность $U$ точки $x_0$ такая, что
$$\forall x \in (U \cap A) \setminus \{x_0\} \implies f(x) \in V.$$
\end{definition}

\begin{definition}[Предел последовательности]
Пусть $(X, \mathcal{T})$ --- топологическое пространство, а $\{x_n\}_{n \in \mathbb{N}}$ --- последовательность точек из $X$.
Точка $x_0 \in X$ называется \textbf{пределом последовательности} $\{x_n\}$, если для любой окрестности $U$ точки $x_0$ существует $m \in \mathbb{N}$ такое, что для всех $n > m$ выполняется $x_n \in U$.
\end{definition}

\begin{lemma}[Предел последовательности в метрике]
Пусть $(X, d)$ --- метрическое пространство, а $\{x_n\}_{n \in \mathbb{N}}$ --- последовательность точек из $X$. Тогда следующие условия эквивалентны:
\begin{enumerate}
    \item $x_0 = \lim x_n$ при $n \to \infty$.
    \item Для любого $\varepsilon > 0$ существует $m \in \mathbb{N}$ такое, что для всех $n > m$ выполнено $x_n \in B(x_0, \varepsilon)$.
    \item Для любого $\varepsilon > 0$ существует $m \in \mathbb{N}$ такое, что для всех $n > m$ выполнено $d(x_0, x_n) < \varepsilon$.
\end{enumerate}
\end{lemma}

\begin{theorem}[О внутренности и точках прикосновения в метрическом пространстве]
Пусть $(X, d)$ --- метрическое пространство, $A \subset X$.
\begin{enumerate}
    \item Точка $a \in A$ является внутренней тогда и только тогда, когда существует $\varepsilon > 0$ такое, что $B(a, \varepsilon) \subset A$.
    \item Для точки $x_0 \in X$ эквивалентны условия:
    \begin{itemize}
        \item $x_0$ --- точка прикосновения $A$.
        \item Для любого $\varepsilon > 0$ существует $a \in A$ такое, что $d(x_0, a) < \varepsilon$.
        \item Существует последовательность $\{a_n\} \subset A$, сходящаяся к $x_0$.
    \end{itemize}
    \item Для изолированной точки $a \in A$ существует $\varepsilon > 0$ такое, что $B(a, \varepsilon) \cap A = \{a\}$.
\end{enumerate}
\end{theorem}

\begin{theorem}[Критерий непрерывности]
Пусть $f: X \to Y$ --- отображение топологических пространств. Тогда эквивалентны следующие условия:
\begin{enumerate}
    \item $f$ непрерывно в каждой точке $X$.
    \item Для любого открытого множества $V \subset Y$ прообраз $f^{-1}(V)$ открыт в $X$.
    \item Для любого замкнутого множества $C \subset Y$ прообраз $f^{-1}(C)$ замкнут в $X$.
\end{enumerate}
\end{theorem}

\begin{theorem}[Критерии непрерывности в метрических пространствах]
Пусть $f: X \to Y$ --- отображение метрических пространств. Тогда эквивалентны:
\begin{enumerate}
    \item $f$ непрерывно в точке $x_0$.
    \item Для любого $\varepsilon > 0$ существует окрестность $U$ точки $x_0$ такая, что для всех $x \in U$ выполнено $d(f(x), f(x_0)) < \varepsilon$.
    \item Для любого $\varepsilon > 0$ существует $\delta > 0$ такое, что из $d(x, x_0) < \delta$ следует $d(f(x), f(x_0)) < \varepsilon$.
\end{enumerate}
\end{theorem}

\subsection{Подпространства, базы и предбазы}

\begin{definition}[Индуцированная топология]
Пусть $(X, \mathcal{T})$ --- топологическое пространство, $A \subset X$. Обозначим через $\mathcal{T}_A$ семейство подмножеств $A$, заданных условием
$$W \in \mathcal{T}_A \iff \exists V \in \mathcal{T} : W = V \cap A.$$
Тогда $\mathcal{T}_A$ является топологией на $A$ и называется \textbf{индуцированной топологией}. Пару $(A, \mathcal{T}_A)$ называют \textbf{подпространством} пространства $(X, \mathcal{T})$.
\end{definition}

\begin{theorem}
Множество $D \subset A$ замкнуто в подпространстве $(A, \mathcal{T}_A)$ тогда и только тогда, когда существует замкнутое множество $C \subset X$ такое, что $D = C \cap A$.
\end{theorem}

\begin{theorem}[Об индуцированной топологии через отображение]
Пусть $(X, \mathcal{T})$ --- топологическое пространство, $Y$ --- множество, и $f: Y \to X$ --- отображение. Определим
$$\mathcal{T}_f = \{V \subset Y \mid \exists U \in \mathcal{T}: V = f^{-1}(U)\}.$$
Тогда $\mathcal{T}_f$ является топологией на $Y$.
\end{theorem}

\begin{lemma}[О замкнутых множествах в индуцированной топологии]
Пусть $\mathcal{T}_f$ --- индуцированная топология из предыдущей теоремы. Тогда множество $B \subset Y$ замкнуто в $(Y, \mathcal{T}_f)$ тогда и только тогда, когда существует замкнутое множество $A \subset X$ такое, что $B = f^{-1}(A)$.
\end{lemma}

\begin{definition}[База и предбаза]
Пусть $(X, \mathcal{T})$ --- топологическое пространство.
\begin{itemize}
    \item Множество $\mathcal{B} \subset \mathcal{T}$ называется \textbf{базой} топологии $\mathcal{T}$, если любой элемент $V \in \mathcal{T}$ представим в виде объединения элементов из $\mathcal{B}$.
    \item Множество открытых подмножеств $\mathcal{P} \subset \mathcal{T}$ называется \textbf{предбазой}, если конечные пересечения элементов $\mathcal{P}$ образуют базу топологии.
\end{itemize}
\end{definition}

\begin{lemma}[О базисе из шаров]
Пусть $(X, d)$ --- произвольное метрическое пространство, а $\mathcal{T}_d$ --- соответствующая метрическая топология. Тогда:
\begin{enumerate}
    \item $\{B(x, r) \mid x \in X,\ r > 0\}$ является базой $\mathcal{T}_d$.
    \item $\{B(x, r) \mid x \in X,\ r \in \mathbb{Q}_{>0}\}$ также является базой $\mathcal{T}_d$.
    \item $\left\{B\left(x, \frac{1}{n}\right) \mid x \in X,\ n \in \mathbb{N}\right\}$ является базой $\mathcal{T}_d$.
\end{enumerate}
\end{lemma}

\section{Лекция 2. Вектор-функции и кривые}

\subsection{Дифференцирование и интегрирование вектор-функций}

\begin{definition}[Производная вектор-функции]
Пусть $\vec{F}: (\alpha, \beta) \to V$ --- функция со значениями в векторном пространстве $V$ со скалярным произведением. Вектор $\vec{a} = \vec{F}'(t_0)$ называется производной вектор-функции в точке $t_0$, если
$$\lim_{t \to t_0} \frac{\vec{F}(t) - \vec{F}(t_0)}{t - t_0} = \vec{a}.$$
\end{definition}

\begin{lemma}[Свойства производных вектор-функций]
Пусть $\vec{f}, \vec{g}: (\alpha, \beta) \to V$ --- дифференцируемые вектор-функции, $u: (\alpha, \beta) \to \mathbb{R}$ --- дифференцируемая скалярная функция. Тогда:
\begin{enumerate}
    \item $(\vec{f} + \vec{g})'(t) = \vec{f}'(t) + \vec{g}'(t)$.
    \item $(c\vec{f})'(t) = c\vec{f}'(t)$ для любого $c \in \mathbb{R}$.
    \item $(\vec{f}, \vec{g})'(t) = (\vec{f}'(t), \vec{g}(t)) + (\vec{f}(t), \vec{g}'(t))$.
    \item $(u\vec{f})'(t) = u'(t)\vec{f}(t) + u(t)\vec{f}'(t)$.
\end{enumerate}
\end{lemma}

\begin{theorem}[Свойства интегралов от вектор-функций]
Пусть $\vec{f}: [\alpha, \beta] \to V$ --- непрерывная вектор-функция, а $\vec{F}$ --- ее первообразная. Тогда:
\begin{enumerate}
    \item $\displaystyle \int_{\alpha}^{\beta} \vec{f}(t)\,dt = \vec{F}(\beta) - \vec{F}(\alpha)$.
    \item $\displaystyle \left\| \int_{\alpha}^{\beta} \vec{f}(t)\,dt \right\| \le \int_{\alpha}^{\beta} \|\vec{f}(t)\|\,dt$.
    \item Если $\vec{f}(t) = f^1(t)\vec{a}_1 + \dots + f^n(t)\vec{a}_n$ в базисе $V$, то интегрирование выполняется покоординатно.
\end{enumerate}
\end{theorem}

\begin{proposition}
Каждая непрерывная вектор-функция $\vec{f}: [\alpha, \beta] \to V$ обладает первообразной. Если $\vec{F}_1$ и $\vec{F}_2$ --- две первообразные одной и той же функции, то $\vec{F}_1(t) - \vec{F}_2(t)$ не зависит от $t$.
\end{proposition}

\subsection{Параметризованные кривые}

Рассматривается евклидово пространство $E = \mathbb{R}^3$ со стандартным началом координат $O$ и радиус-векторами точек.

\begin{definition}[Параметризованная кривая]
Непрерывное отображение $\gamma: [\alpha, \beta] \to E$ называется \textbf{параметризованной кривой}. Точка $A = \gamma(\alpha)$ называется началом кривой, а точка $B = \gamma(\beta)$ --- ее концом.
\end{definition}

\begin{definition}[Касательная прямая]
Прямая $l$ называется \textbf{касательной} к кривой $\gamma(t)$ в точке $t_0$, если
$$\lim_{t \to t_0} \frac{h(t)}{\|\gamma(t) - \gamma(t_0)\|} = 0,$$
где $h(t)$ --- кратчайшее расстояние от точки $\gamma(t)$ до прямой $l$.
\end{definition}

\begin{theorem}[Касательная прямая]
Если кривая задана дифференцируемой вектор-функцией
$\gamma(t) = (x(t), y(t), z(t))$, то касательная прямая
в точке $A = \gamma(t_0)$ имеет параметрический вид
$$
(X, Y, Z) = (x(t_0), y(t_0), z(t_0)) + \lambda\,(x'(t_0), y'(t_0), z'(t_0)).
$$
Если соответствующие производные не равны нулю, это можно переписать и в каноническом виде.
\end{theorem}

\begin{definition}[Длина кривой]
Длина спрямляемой дифференцируемой кривой на промежутке $[\alpha, \beta]$ вычисляется по формуле
$$L(\gamma) = \int_{\alpha}^{\beta} \|\gamma'(t)\|\,dt = \int_{\alpha}^{\beta} \sqrt{(x'(t))^2 + (y'(t))^2 + (z'(t))^2}\,dt.$$
\end{definition}

\subsection{Естественная параметризация и формулы Френе}

\begin{definition}[Естественная параметризация]
Пусть $\gamma(t)$ --- кривая, заданная параметрически. Тогда $\tilde{\gamma}(s)$ называется \textbf{естественной параметризацией} кривой, если модуль ее скорости равен $1$:
$$\left\| \frac{d\tilde{\gamma}}{ds} \right\| = 1.$$
\end{definition}

\begin{proposition}[О свойствах естественной параметризации]
Пусть $\tilde{\gamma}$ --- естественная параметризация кривой $\gamma$. Тогда:
\begin{enumerate}
    \item $\dfrac{d\gamma}{dt} = \dfrac{d\tilde{\gamma}}{ds}\cdot \dfrac{ds}{dt}$.
    \item Если $A = \gamma(s_0)$, $B = \gamma(s_1)$, то длина дуги между $A$ и $B$ равна $s_1 - s_0$.
\end{enumerate}
\end{proposition}

\begin{theorem}[О существовании естественной параметризации]
Пусть $\vec{\gamma}(t)$ --- регулярная кривая, то есть $\vec{\gamma}$ непрерывно дифференцируема и $\vec{\gamma}'(t) \neq 0$ для всех $t$. Тогда существует естественный параметр кривой $\vec{\gamma}$.
\end{theorem}

\begin{proof}[Идея доказательства]
Определим
$$s(t) = \int_{\alpha}^{t} \|\gamma'(\tau)\|\,d\tau.$$
Поскольку $\|\gamma'(t)\| > 0$, функция $s(t)$ строго возрастает и потому обратима. Обратная функция позволяет перепараметризовать кривую по длине дуги.
\end{proof}

\begin{definition}[Кривизна, нормаль и бинормаль]
Пусть $\tilde{\gamma}$ --- естественная параметризация кривой. Обозначим
$$\vec{\tau} = \frac{d\tilde{\gamma}}{ds}.$$
Тогда $\vec{\tau}$ называется единичным касательным вектором. \textbf{Кривизной} кривой называется
$$\kappa = \left\| \frac{d\vec{\tau}}{ds} \right\|.$$
Если $\kappa \neq 0$, то единичный вектор главной нормали определяется равенством
$$\frac{d\vec{\tau}}{ds} = \kappa \vec{n},$$
а вектор
$$\vec{b} = \vec{\tau} \times \vec{n}$$
называется \textbf{бинормалью}.
\end{definition}

\begin{theorem}[Формулы Френе]
Пусть $\gamma$ --- регулярная кривая, а $\tilde{\gamma}$ --- ее естественная параметризация. Тогда для базиса Френе $\{\vec{\tau}, \vec{n}, \vec{b}\}$ выполняются формулы
\begin{align*}
    \frac{d\vec{\tau}}{ds} &= \kappa \vec{n}, \\
    \frac{d\vec{n}}{ds} &= -\kappa \vec{\tau} + \varkappa \vec{b}, \\
    \frac{d\vec{b}}{ds} &= -\varkappa \vec{n},
\end{align*}
где $\varkappa$ --- кручение.
\end{theorem}

\begin{proposition}
Пусть $\gamma$ --- регулярная непрерывно дифференцируемая кривая.
\begin{enumerate}
    \item Если $\kappa \equiv 0$, то все точки кривой лежат на некоторой прямой.
    \item Если $\varkappa \equiv 0$, то кривая лежит в некоторой плоскости.
\end{enumerate}
\end{proposition}

\begin{example}[Окружность]
Для кривой $\gamma(t) = (R \cos t, R \sin t, 0)$ имеем
$$\gamma'(t) = (-R \sin t, R \cos t, 0), \qquad \|\gamma'(t)\| = R.$$
Ее естественная параметризация задается формулой
$$\tilde{\gamma}(s) = \left(R \cos \frac{s}{R}, R \sin \frac{s}{R}, 0\right),$$
а кривизна равна $\kappa = 1/R$.
\end{example}

\begin{example}[Винтовая линия]
Для кривой $\gamma(t) = (a \cos t, a \sin t, bt)$ получаем
$$\|\gamma'(t)\| = \sqrt{a^2 + b^2},$$
а естественная параметризация имеет вид
$$
\tilde{\gamma}(s) = \left(a \cos \frac{s}{\sqrt{a^2+b^2}},\ a \sin \frac{s}{\sqrt{a^2+b^2}},\ \frac{bs}{\sqrt{a^2+b^2}}\right).
$$
\end{example}

\section{Лекция 3. Билинейные и квадратичные формы}

\subsection{Билинейные формы}

\begin{definition}[Билинейная форма]
Отображение $h: L \times L \to \mathbb{R}$ называется \textbf{билинейной формой}, если оно линейно по каждому аргументу:
\begin{enumerate}
    \item $h(x+y, z) = h(x, z) + h(y, z)$ и $h(x, y+z) = h(x, y) + h(x, z)$.
    \item $h(\alpha x, y) = \alpha h(x, y) = h(x, \alpha y)$.
\end{enumerate}
\end{definition}

\begin{definition}[Матрица билинейной формы]
Пусть $(a) = \{a_1, \dots, a_n\}$ --- базис пространства $L$. Матрицей билинейной формы в базисе $(a)$ называется матрица $[h]_{(a)} = (h_{ij})$, где $h_{ij} = h(a_i, a_j)$.
\end{definition}

\begin{definition}[Невырожденность и симметрия]
Билинейная форма $g: L \times L \to \mathbb{R}$ называется:
\begin{itemize}
    \item \textbf{невырожденной}, если из $g(a, x) = 0$ для всех $x \in L$ следует $a = 0$;
    \item \textbf{симметричной}, если $g(x, y) = g(y, x)$ для всех $x, y \in L$.
\end{itemize}
\end{definition}

\begin{lemma}[Свойства билинейных форм]
Пусть $h, g: L \times L \to \mathbb{R}$ --- билинейные формы, а $(a)$ --- базис пространства $L$. Тогда:
\begin{enumerate}
    \item Для любых $x,y \in L$ выполняется
    $$h(x,y) = [x]_{(a)} [h]_{(a)} [y]_{(a)}^T = \sum_{i=1}^n \sum_{j=1}^n h_{ij} x^i y^j.$$
    \item $h$ симметрична тогда и только тогда, когда матрица $[h]_{(a)}$ симметрична.
    \item Если $C$ --- матрица перехода от базиса $(a)$ к базису $(b)$, то
    $$[h]_{(b)} = C [h]_{(a)} C^T.$$
    \item Для координат векторов выполняется правило $[x]_{(a)} = [x]_{(b)} C$.
    \item Если $h(x,y) = g(Bx, y)$ для линейного оператора $B$, то $[h]_{(a)} = [B]_{(a)} [g]_{(a)}$.
    \item Форма $g$ невырождена тогда и только тогда, когда $\det[g]_{(a)} \neq 0$.
    \item Для линейного оператора $B$ выполняется $[B(x)]_{(a)} = [x]_{(a)} [B]_{(a)}$.
\end{enumerate}
\end{lemma}

\begin{lemma}
Пусть $g: L \times L \to \mathbb{R}$ --- невырожденная билинейная форма. Тогда для любого линейного функционала $f: L \to \mathbb{R}$ существует единственный вектор $a_f \in L$ такой, что
$$f(y) = g(a_f, y) \quad \forall y \in L.$$
\end{lemma}

\begin{theorem}[Представление билинейной формы оператором]
Пусть $h, g: L \times L \to \mathbb{R}$ --- билинейные формы, причем $g$ невырождена. Тогда существует единственный линейный оператор $B: L \to L$ такой, что
$$h(x, y) = g(Bx, y).$$
Если $h$ симметрична, то оператор $B$ самосопряжен относительно $g$.
\end{theorem}

\subsection{Квадратичные формы}

\begin{definition}[Квадратичная форма]
Квадратичной формой от $n$ переменных над полем $k$ называется однородный многочлен степени $2$:
$$Q(x^1, \dots, x^n) = \sum_{i=1}^n q_{ii}(x^i)^2 + 2\sum_{i<j} q_{ij}x^i x^j,$$
где $q_{ij} \in k$.
\end{definition}

\begin{example}
Для квадратичной формы
$$Q(x, y, z) = 3x^2 - 4xy + 5xz + 2z^2 - 3xz$$
соответствующая матрица имеет вид
$$[Q] = \begin{pmatrix} 3 & -2 & 1 \\ -2 & 0 & 0 \\ 1 & 0 & 2 \end{pmatrix}.$$
\end{example}

\begin{theorem}[Связь билинейной и квадратичной форм]
\begin{enumerate}
    \item Если $h: L \times L \to k$ --- билинейная форма, то $Q(x) = h(x, x)$ является квадратичной формой.
    \item Для любой квадратичной формы $Q$ существует единственная симметричная билинейная форма $h$ такая, что $Q(x) = h(x, x)$.
    \item Эта билинейная форма выражается формулой поляризации:
    $$h(x, y) = \frac{1}{2}\big(Q(x+y) - Q(x) - Q(y)\big).$$
\end{enumerate}
\end{theorem}

\begin{proposition}[Симметризация]
Если $h$ --- билинейная форма, то
$$\tilde{h}(x, y) = \frac{1}{2}\big(h(x, y) + h(y, x)\big)$$
тоже является билинейной формой, причем $\tilde{h}(x, x) = h(x, x)$.
\end{proposition}

\section{Лекция 4. Поверхности в \texorpdfstring{$\mathbb{R}^3$}{R3}}

\subsection{График функции как поверхность}

\begin{definition}[Элементарная область]
Открытое множество $V \subset \mathbb{R}^2$ называется \textbf{элементарной областью}, если оно гомеоморфно единичному открытому кругу
$$B^2 = \{(x,y) \in \mathbb{R}^2 \mid x^2 + y^2 < 1\}.$$
\end{definition}

\begin{definition}[Гладкая поверхность]
Множество $S \subset \mathbb{R}^3$ называется \textbf{гладкой поверхностью}, если для любой точки $p \in S$ существует окрестность $W \subset \mathbb{R}^3$ и гомеоморфизм $r: V \to S \cap W$, где $V$ --- элементарная область в $\mathbb{R}^2$, такой что:
\begin{enumerate}
    \item $r(u, v)$ гладкая;
    \item $\dfrac{\partial r}{\partial u} \times \dfrac{\partial r}{\partial v} \neq \mathbf{0}$;
    \item матрица Якоби имеет ранг $2$.
\end{enumerate}
\end{definition}

\begin{theorem}[График функции как поверхность]
Пусть $V \subset \mathbb{R}^2$ --- открытое множество, $f: V \to \mathbb{R}$ --- непрерывная функция. Тогда график
$$\Gamma_f = \{(u, v, f(u,v)) \mid (u,v) \in V\}$$
является поверхностью. Более того, отображение
$$r(u,v) = (u, v, f(u,v))$$
задает параметризацию $\Gamma_f$ в окрестности каждой ее точки.
\end{theorem}

\begin{proof}[Идея доказательства]
Рассматриваем проекцию $p: \mathbb{R}^3 \to \mathbb{R}^2$, $p(x^1, x^2, x^3) = (x^1, x^2)$, и ее сужение на график. Оно является обратным к отображению $r$, поэтому график гомеоморфен $V$. Если $f$ достаточно гладкая, то параметризация регулярна.
\end{proof}

\subsection{Касательное пространство и касательная плоскость}

\begin{definition}[Касательный вектор]
Пусть $S \subset \mathbb{R}^3$ --- поверхность, $P \in S$. Вектор $\vec{a}$ называется \textbf{касательным вектором} к $S$ в точке $P$, если существует кривая $\gamma$, лежащая в $S$, такая что $\gamma(t_0) = P$ и $\vec{a} = \gamma'(t_0)$.
\end{definition}

\begin{definition}[Касательное пространство]
Множество всех касательных векторов к $S$ в точке $P$ называется \textbf{касательным пространством} и обозначается $T_P S$.
\end{definition}

\begin{theorem}[Базис касательного пространства]
Пусть $S \subset \mathbb{R}^3$ --- поверхность, $r: V \to \mathbb{R}^3$ --- регулярная локальная параметризация поверхности, $P = r(u_0, v_0) \in S$. Тогда $T_P S$ является двумерным векторным пространством, а векторы
$$r_u = \frac{\partial r}{\partial u}(u_0, v_0), \qquad r_v = \frac{\partial r}{\partial v}(u_0, v_0)$$
образуют его базис.
\end{theorem}

\begin{lemma}[О задании кривых на поверхности]
Пусть $r: V \to \mathbb{R}^3$ --- регулярная параметризация поверхности $S$, $P = r(u_0, v_0)$, а $\gamma$ --- кривая на $S$ такая, что $\gamma(t_0) = P$. Тогда в некоторой окрестности точки $t_0$ кривая представима в виде
$$\gamma(t) = r(u(t), v(t)).$$
\end{lemma}

\begin{definition}[Касательная плоскость]
Касательной плоскостью к регулярной поверхности в точке $P$ называется множество всех точек $X$, для которых вектор $PX$ принадлежит $T_P S$.
\end{definition}

\begin{proposition}[Уравнение касательной плоскости]
Пусть $S \subset \mathbb{R}^3$ --- поверхность, $r: V \to \mathbb{R}^3$ --- ее локальная регулярная параметризация, $r(u_0, v_0)=P=(x_0,y_0,z_0)$. Тогда касательная плоскость к $S$ в точке $P$ задается уравнением
$$
\det \begin{pmatrix}
x-x_0 & y-y_0 & z-z_0 \\
\frac{\partial r^1}{\partial u} & \frac{\partial r^2}{\partial u} & \frac{\partial r^3}{\partial u} \\
\frac{\partial r^1}{\partial v} & \frac{\partial r^2}{\partial v} & \frac{\partial r^3}{\partial v}
\end{pmatrix} = 0.
$$
\end{proposition}

\begin{definition}[Нормаль]
Вектор $\vec{m}$, перпендикулярный касательной плоскости, называется единичным вектором нормали к поверхности $S$ в точке $P$. Прямая, проходящая через $P$ параллельно $\vec{m}$, называется нормалью к поверхности.
\end{definition}

\subsection{Поверхности вращения}

\begin{lemma}[О поверхности вращения]
Пусть $\gamma(t) = (f(t), 0, g(t))$. Тогда поверхность вращения этой кривой вокруг оси $Z$ задается параметризацией
$$r(u, v) = \big(f(u)\cos v,\ f(u)\sin v,\ g(u)\big).$$
\end{lemma}

\begin{example}[Сфера]
Если $\gamma(t) = (R \cos t, 0, R \sin t)$, то
$$r(u, v) = \big(R \cos u \cos v,\ R \cos u \sin v,\ R \sin u\big)$$
является параметрическим уравнением сферы.
\end{example}

\begin{example}[Тор]
Если $\gamma(t) = (R + r\cos t, 0, r\sin t)$, то
$$r(u, v) = \big((R + r\cos u)\cos v,\ (R + r\cos u)\sin v,\ r\sin u\big)$$
задает тор.
\end{example}

\subsection{Первая и вторая фундаментальные формы}

\begin{definition}[Первая фундаментальная форма]
Пусть $S$ --- поверхность, $p \in S$, $T_pS$ --- касательное пространство в точке $p$. \textbf{Первая основная форма} поверхности $S$ в точке $p$ определяется как
$$I_p(\bar{a}, \bar{b}) = \langle \bar{a}, \bar{b} \rangle, \qquad \bar{a}, \bar{b} \in T_pS.$$
\end{definition}

\begin{proposition}[Коэффициенты первой фундаментальной формы]
Пусть $r: V \to \mathbb{R}^3$ --- регулярная локальная параметризация поверхности $S$, $p = r(u_0, v_0)$. Тогда векторы $r_u$ и $r_v$ образуют базис $T_pS$, а коэффициенты первой фундаментальной формы равны
$$g_{11} = \langle r_u, r_u \rangle,\qquad g_{12} = g_{21} = \langle r_u, r_v \rangle,\qquad g_{22} = \langle r_v, r_v \rangle.$$
Если $\bar{a} = a^1 r_u + a^2 r_v$, то
$$I_p(\bar{a}) = g_{11}(a^1)^2 + 2g_{12}a^1a^2 + g_{22}(a^2)^2.$$
\end{proposition}

\begin{definition}[Вторая фундаментальная форма]
Пусть $S$ --- поверхность, $r$ --- регулярная параметризация, $r(u_0, v_0) = p \in S$. Для вектора $\bar{a} = a^1 r_u + a^2 r_v \in T_pS$ определим
$$II_p(\bar{a}) = b_{11}(a^1)^2 + 2b_{12}a^1a^2 + b_{22}(a^2)^2,$$
где
$$b_{11} = \langle r_{uu}, m \rangle,\qquad b_{12} = \langle r_{uv}, m \rangle,\qquad b_{22} = \langle r_{vv}, m \rangle,$$
а $m$ --- единичная нормаль к поверхности.
\end{definition}

\begin{lemma}[Проекция ускорения]
Пусть $\gamma(t)$ --- кривая на поверхности $S$, $\gamma(t_0)=p$,
причем $\gamma(t)=r(u(t),v(t))$. Тогда
$$
\langle \gamma''(t_0), m \rangle = b_{11}(u'(t_0))^2 + 2b_{12}u'(t_0)v'(t_0) + b_{22}(v'(t_0))^2.
$$
\end{lemma}

\begin{lemma}[О знаке второй фундаментальной формы]
Вторая основная форма поверхности в точке $p$ при замене параметров либо не меняется, либо умножается на $-1$. Знак меняется только при замене ориентации нормали $m \mapsto -m$.
\end{lemma}

\begin{theorem}[Теорема Менье]
Пусть $\gamma(t) \in S$ --- кривая на поверхности $S$, проходящая через точку $p$. Тогда нормальная кривизна $k_n$ выражается как
$$k_n = \frac{II_p(\gamma'(t_0))}{I_p(\gamma'(t_0))}.$$
\end{theorem}

\begin{theorem}[Теорема Эйлера]
Пусть $S$ --- регулярная поверхность, $p \in S$. Тогда в касательной плоскости существуют два взаимно перпендикулярных направления $l_1$ и $l_2$, в которых нормальная кривизна принимает экстремальные значения $k_1$ и $k_2$. Для любого направления $l$, образующего угол $\theta$ с главным направлением $l_1$, справедлива формула
$$k_l = k_1 \cos^2 \theta + k_2 \sin^2 \theta.$$
\end{theorem}

\subsection{Деривационные уравнения и уравнения Гаусса-Кодацци}

\begin{proposition}[Уравнения Гаусса-Вейнгартена]
Деривационные уравнения поверхности можно записать в матричной форме:
$$
\frac{\partial}{\partial u}
\begin{pmatrix}
\bar{r}_u\\
\bar{r}_v\\
\bar{m}
\end{pmatrix}
=
\begin{pmatrix}
\Gamma_{11}^1 & \Gamma_{11}^2 & b_{11}\\
\Gamma_{12}^1 & \Gamma_{12}^2 & b_{12}\\
-b_1^1 & -b_1^2 & 0
\end{pmatrix}
\begin{pmatrix}
\bar{r}_u\\
\bar{r}_v\\
\bar{m}
\end{pmatrix},
$$
$$
\frac{\partial}{\partial v}
\begin{pmatrix}
\bar{r}_u\\
\bar{r}_v\\
\bar{m}
\end{pmatrix}
=
\begin{pmatrix}
\Gamma_{12}^1 & \Gamma_{12}^2 & b_{12}\\
\Gamma_{22}^1 & \Gamma_{22}^2 & b_{22}\\
-b_2^1 & -b_2^2 & 0
\end{pmatrix}
\begin{pmatrix}
\bar{r}_u\\
\bar{r}_v\\
\bar{m}
\end{pmatrix}.
$$
\end{proposition}

\begin{theorem}[Об уравнениях Гаусса-Кодацци]
Имеют место уравнения согласованности
$$\frac{\partial A_2}{\partial u} - \frac{\partial A_1}{\partial v} + [A_1, A_2] = 0,$$
где $A_1$ и $A_2$ --- матрицы связности деривационных уравнений.
\end{theorem}

\begin{theorem}[Теорема Гаусса]
Гауссова кривизна выражается только через компоненты первой фундаментальной формы $g_{ij}$ и их частные производные:
$$K_g = \frac{b_{11}b_{22} - b_{12}^2}{g_{11}g_{22} - g_{12}^2} = \frac{\det(II)}{\det(I)}.$$
\end{theorem}

\end{document}
